\documentclass{article}
\usepackage[utf8]{inputenc}
\usepackage[T1]{fontenc}
\usepackage{amsmath,amssymb}

\begin{document}

\section*{Constantes}
\begin{itemize}
    \item Fontes \(S_d\): especialistas que passam conhecimento em \(d\);
    \item Demandantes \(D_d\): quem precisa de conhecimento em \(d\);
    \item Possivelmente peso \(w(v,d)\): criticidade da disciplina para o projeto do nó \(v\).
\end{itemize}

\section*{Variáveis}
\[
    x_e \in \{0,1\}, \quad e \in C,
\]
\[
    r(v,d) \in \{0,1\}: \text{\(v\) é alcançado em \(d\)}.
\]

\section*{Objetivo}
\[
    \max \sum_d \sum_v w(v,d) \cdot r(v,d)
\]

\section*{Restrições}
\[
    \sum_{e \in C} x_e \le B
\]
\[
    r(v,d) = 1 \quad \text{para } v \in S_d.
\]

Determinar se \(v\) é alcançado em \(d\) se existe caminho até \(v\) usando possivelmente limite de saltos; possivelmente máximo de conhecimento que sai de um especialista para não ligar tudo à mesma pessoa.

\end{document}